% =====================================================================
%  DC-FOUNDATION-01 — Raspberry Pi Node Foundation
%  Lab Execution Record & Meta Report  (P00 / P01 · BSA-conformant)
%  Author: Redji Jean Baptiste (Toussaint)
%  Compile:  pdflatex DC-FOUNDATION-01.tex   (run twice for ToC/refs)
%
%  IMAGE NOTE ---------------------------------------------------------
%  Figures are referenced by their mapping-table names (fig-*.png).
%  Rename the screenshots to those names and place them in ./evidence/
%  (see the mapping table at the end of the source, or Appendix C).
%  Any missing image will show a gray placeholder box, not an error,
%  because we use \IfFileExists guards below.
% =====================================================================
\documentclass[11pt,a4paper]{article}

\usepackage[T1]{fontenc}
\usepackage[utf8]{inputenc}
\usepackage{lmodern}
\usepackage{amssymb}
\usepackage[a4paper,margin=1in,headheight=14pt]{geometry}
\usepackage{graphicx}
\usepackage{xcolor}
\usepackage{booktabs}
\usepackage{longtable}
\usepackage{array}
\usepackage{tabularx}
\usepackage{enumitem}
\usepackage{listings}
\usepackage{fancyhdr}
\usepackage{titlesec}
\usepackage{caption}
\usepackage{float}
\usepackage[hidelinks,colorlinks=true,linkcolor=dcred,urlcolor=dcblue,citecolor=dcblue]{hyperref}

% ---- palette --------------------------------------------------------
\definecolor{dcred}{RGB}{155,25,45}
\definecolor{dcblue}{RGB}{20,80,140}
\definecolor{dcgreen}{RGB}{20,120,60}
\definecolor{dcgray}{RGB}{90,90,90}
\definecolor{codebg}{RGB}{245,245,245}
\definecolor{consolebg}{RGB}{25,30,40}
\definecolor{consolefg}{RGB}{225,225,225}
\definecolor{okgreen}{RGB}{20,120,60}
\definecolor{failred}{RGB}{170,30,30}
\definecolor{warnamber}{RGB}{170,110,10}

% ---- section styling ------------------------------------------------
\titleformat{\section}{\Large\bfseries\color{dcred}}{\thesection}{0.6em}{}
\titleformat{\subsection}{\large\bfseries\color{dcblue}}{\thesubsection}{0.6em}{}
\titleformat{\subsubsection}{\normalsize\bfseries\color{dcgray}}{\thesubsubsection}{0.6em}{}

% ---- headers --------------------------------------------------------
\pagestyle{fancy}
\fancyhf{}
\lhead{\small\textcolor{dcgray}{DC-FOUNDATION-01 · dc-node-01}}
\rhead{\small\textcolor{dcgray}{\thepage}}
\lfoot{\small\textcolor{dcgray}{BSA / P00 Foundation}}
\rfoot{\small\textcolor{dcgray}{Redji Jean Baptiste}}
\renewcommand{\headrulewidth}{0.4pt}

% ---- listings styles ------------------------------------------------
\lstdefinestyle{bash}{
  backgroundcolor=\color{codebg},
  basicstyle=\ttfamily\small,
  keywordstyle=\color{dcblue}\bfseries,
  commentstyle=\color{dcgreen}\itshape,
  breaklines=true, columns=fullflexible,
  frame=single, rulecolor=\color{dcgray!40},
  showstringspaces=false, xleftmargin=6pt, xrightmargin=6pt,
  aboveskip=8pt, belowskip=8pt,
  morekeywords={sudo,nmcli,ssh,scp,systemctl,apt,ufw,git,docker,timedatectl,chmod,cat,grep,ls}
}
\lstdefinestyle{console}{
  backgroundcolor=\color{consolebg},
  basicstyle=\ttfamily\footnotesize\color{consolefg},
  breaklines=true, columns=fullflexible,
  frame=single, rulecolor=\color{consolebg},
  showstringspaces=false, xleftmargin=6pt, xrightmargin=6pt,
  aboveskip=8pt, belowskip=8pt
}
\lstset{style=bash}

% ---- helper macros --------------------------------------------------
\graphicspath{{evidence/}{./}}

% safe figure: shows a placeholder if the file is missing
\newcommand{\evfig}[3]{% {filename}{width-fraction}{caption}
  \begin{figure}[H]\centering
  \IfFileExists{evidence/#1}{\includegraphics[width=#2\linewidth]{#1}}%
  {\fbox{\parbox[c][3.2cm][c]{#2\linewidth}{\centering\ttfamily\small\textcolor{dcgray}{[ figure placeholder — supply \detokenize{#1} ]}}}}%
  \caption{#3}\end{figure}}

\newcommand{\pass}{\textcolor{okgreen}{\textbf{$\checkmark$ PASS}}}
\newcommand{\fail}{\textcolor{failred}{\textbf{$\times$ FAIL}}}
\newcommand{\open}{\textcolor{warnamber}{\textbf{OPEN}}}
\newcommand{\closed}{\textcolor{okgreen}{\textbf{CLOSED}}}
\newcommand{\bsl}[1]{\textcolor{dcred}{\textbf{BSL-#1}}}
\newcommand{\domain}[1]{\textsc{#1}}

\newcommand{\parttitle}[2]{%
  \clearpage
  \begin{center}
  {\Huge\bfseries\color{dcred} #1}\\[4pt]
  {\large\color{dcgray} #2}
  \end{center}
  \vspace{4pt}\hrule\vspace{10pt}}

% =====================================================================
\begin{document}

% ---------------------------------------------------------------- TITLE
\begin{titlepage}
\centering
\vspace*{1.2cm}
{\LARGE\bfseries\color{dcred} DC-FOUNDATION-01}\\[6pt]
{\Large Raspberry Pi Node Foundation}\\[4pt]
{\large\itshape Lab Execution Record \& Meta Report}\\[18pt]
\rule{0.7\linewidth}{0.8pt}\\[14pt]

\begin{tabular}{rl}
\textbf{Project ID}      & P00 (Foundation) — first executable of \textbf{P01, Jan–Apr 2026}\\
\textbf{Node}            & \texttt{dc-node-01} \quad(D-Central node 01)\\
\textbf{Platform}        & Raspberry Pi 5 Model B Rev 1.1 (8\,GB)\\
\textbf{Operating system}& Raspberry Pi OS 64-bit — Debian 13 ``Trixie''\\
\textbf{Kernel}          & 6.18.34+rpt-rpi-2712 (\texttt{arm64})\\
\textbf{Author}          & Redji Jean Baptiste (Toussaint) — 041-269-022\\
\textbf{Standard}        & Master Timeline V22 — Black Start Augmentation System\\
\textbf{STDS layer}      & 0 — Userland (activates P01)\\
\textbf{Build window}    & 2026-07-09 $\rightarrow$ 2026-07-22\\
\end{tabular}\\[20pt]

\fbox{\parbox{0.86\linewidth}{\centering
\textbf{\color{okgreen} STATUS: COMPLETE.} Part A steps 0–12 executed, evidenced,
and verified. SSH hardening proven by negative test. Reboot persistence confirmed.
All nine defects (D1–D9) closed. Published to \texttt{RedjiJB/cst-portfolio}.}}\\[16pt]

{\small\color{dcgray} Three-Track status: CL \pass\ $\rightarrow$ SL \pass\ $\rightarrow$ BSL \pass\ (rebuild demonstrated).
Maturity attained: \bsl{1}, touching \bsl{2}.}

\vfill
{\small\color{dcgray} Companion research seed \textbf{T1.1} — \emph{Minimum Viable Sovereign
Infrastructure Specification} $\rightarrow$ repo \texttt{Sovereign-Bootstrap-v1}.}
\end{titlepage}

\tableofcontents
\clearpage

% ================================================================= §0
\section{Timeline Alignment — Conformance to the BSA Standard}
This report is the first executed artifact of the Master Timeline V22. It is
therefore scored not only for technical correctness but against the
\textbf{Black Start Augmentation System (BSA)} documentation standard: the
Three-Track Rule, the six Operational Domains, BSL maturity, the four
Operational Modes, and the seven-stage Skill Pipeline. This section maps the
P00 work onto that framework so the remainder of the document can be read as
evidence for each claim.

\subsection{The Three-Track Rule}
No skill counts as complete until it clears all three tracks.

\begin{longtable}{@{}p{2.4cm} p{8.4cm} c@{}}
\toprule
\textbf{Track} & \textbf{Requirement / evidence in this build} & \textbf{Status}\\
\midrule\endhead
\textbf{CL — Cognitive} & Theory of hardening, addressing, and container isolation is
explained throughout Parts A/B — e.g.\ the \texttt{sshd\_config} include-ordering
rule (§\ref{sec:d3fix}) and DHCP-vs-static lease semantics (§\ref{sec:baseline}). & \pass\\
\addlinespace
\textbf{SL — System} & A working node exists: hardened SSH, host firewall,
intrusion protection, automated patching, and a container substrate, all
verified live. & \pass\\
\addlinespace
\textbf{BSL — Black Start} & Rebuild-from-zero was not hypothetical: a self-inflicted
network lockout (§\ref{sec:lockout}) forced a full re-image, and the node was
reconstituted from the documented procedure. The changed Machine ID is the proof
(§\ref{sec:rebuild}). & \pass\\
\bottomrule
\end{longtable}

\subsection{Operational Domains — P01 target: score 1 in each}
\begin{longtable}{@{}p{2.3cm} p{7.6cm} c@{}}
\toprule
\textbf{Domain} & \textbf{P00 contribution} & \textbf{Score}\\
\midrule\endhead
\domain{Energy}     & Out of scope for P00; UPS/solar deferred to later P01 tasks. & 0$\rightarrow$pending\\
\domain{Compute}    & Docker substrate operational; \texttt{arm64v8} image path proven (Fig.\ 43). & 1\\
\domain{Storage}    & Portfolio vault on Git; artifacts committed; microSD boot (endurance noted §\ref{sec:storage}). & 1\\
\domain{Network}    & Static addressing achieved; DNS corrected; firewall default-deny. & 1\\
\domain{Control}    & Key-only SSH proven by negative test; least-privilege user model. & 1\\
\domain{Production} & Reproducible build proven by rebuild; Ansible-style repeatability is the next step. & 1\\
\bottomrule
\end{longtable}
\noindent Five of six domains reach the P01 baseline of 1. \domain{Energy} is the
single deferred item and is scheduled with the Kill-a-watt power baseline and
2-node cluster later in P01.

\subsection{BSL Maturity attained}
Phase target for P01–P04 is \bsl{1}/\bsl{2}. This build attains \bsl{1}
(lab-capable) across all in-scope domains and demonstrates elements of \bsl{2}
(offline/recoverable): the node was recovered from bare storage using
documentation alone.

\subsection{Four Operational Modes}
\begin{longtable}{@{}p{3.2cm} p{8.4cm} c@{}}
\toprule
\textbf{Mode} & \textbf{Demonstrated?} & \textbf{Status}\\
\midrule\endhead
1 — Online       & Normal operation over Wi-Fi / Ethernet. & \pass\\
2 — Degraded     & Node traversed three networks (§\ref{sec:relocate}); operated on a phone hotspot. & partial\\
3 — Offline      & Not yet — updates still require internet; local mirror is a later task. & pending\\
4 — Black Start  & The lockout $\rightarrow$ re-image $\rightarrow$ reconstitution is a partial Mode-4 event (no net, storage rebuilt). & partial\\
\bottomrule
\end{longtable}

\subsection{Seven-Stage Skill Pipeline}
\begin{center}\small
\begin{tabular}{@{}c c c c c c c@{}}
\toprule
LEARN & BUILD & INSPECT & BREAK & DEFEND & RECOVER & OPTIMIZE\\
\midrule
\pass & \pass & \pass & \pass & \pass & \pass & partial\\
\bottomrule
\end{tabular}
\end{center}
\noindent \textsc{Inspect} = baseline capture and the \texttt{sshd -T} audit;
\textsc{Break} = the wrong-subnet lockout and the negative auth test;
\textsc{Recover} = the full rebuild. All seven stages are represented, with
\textsc{Optimize} carried forward (SSD migration, DHCP reservation).

\subsection{Black Start Score Matrix — current node state}
\begin{longtable}{@{}p{2.2cm} p{5.4cm} p{4.2cm}@{}}
\toprule
\textbf{Domain} & \textbf{Current state} & \textbf{Target (100)}\\
\midrule\endhead
Energy     & Grid only (0)                          & 72\,hr battery+solar\\
Compute    & Docker local; rebuild $\sim$1 image pass & $<$2\,hr full rebuild\\
Storage    & Git vault + microSD (no M-DISC yet)     & 3-2-1 local + M-DISC\\
Network    & Static + firewall; single path          & 3 independent paths\\
Control    & Key-only SSH, local user; no cloud IdP  & zero cloud dependency\\
Production & Reproducible ($\sim$1 image pass)        & $<$4\,hr from blank HW\\
\bottomrule
\end{longtable}
\noindent \textbf{Rule satisfied:} hard external dependencies were reduced —
password authentication and any cloud-IdP path were removed; the node authenticates
locally by key with no external identity provider.

\subsection{Research seed}
This document is the deliverable for research seed \textbf{T1.1 — Minimum Viable
Sovereign Infrastructure Specification}, and the \texttt{cst-portfolio} repository
is the concrete instance of \texttt{Sovereign-Bootstrap-v1}. The node hostname
\texttt{dc-node-01} designates it as the first node of \textbf{D-Central}.

% ============================================================= PART 0
\parttitle{Part 0}{Image Preparation and First Contact}

\subsection{Tooling}
\begin{center}
\begin{tabular}{@{}lll@{}}
\toprule
\textbf{Tool} & \textbf{Version} & \textbf{Role}\\
\midrule
Raspberry Pi Imager & v2.0.10 (Win 11) & Writes image + first-boot customisation\\
PuTTY               & current          & Initial SSH client\\
OpenSSH (Windows)   & built-in         & \texttt{ssh}/\texttt{scp}/\texttt{ssh-keygen} from PowerShell\\
\bottomrule
\end{tabular}
\end{center}

\subsection{Device, OS, and storage selection}
Raspberry Pi 5 was selected (BCM2712 boot chain). The recommended
\textbf{Raspberry Pi OS 64-bit (Trixie)} image was used — a deliberate deviation
from the manual's Bookworm target, since Bookworm is now the \emph{Legacy} entry
and Trixie receives current security updates (analysed in §\ref{sec:trixie}).

\evfig{fig-0.2-imager-os-trixie64.png}{0.9}{Imager — Raspberry Pi OS 64-bit (Trixie), released 2026-06-18, cached locally.}

The only eligible target was a 59.5\,GB ``USB Mass Storage Device.'' A Windows
AutoPlay prompt claimed the drive was damaged; this is a false positive —
Windows cannot read the Linux \texttt{ext4} partition and must be ignored
(never run the repair tool on a fresh image).

\evfig{fig-0.3-imager-storage-usb.png}{0.9}{Storage selection; note the AutoPlay false-positive warning at lower right.}

\subsection{First-boot customisation}
Hostname \texttt{dc-node-01}, locale Ottawa / America\_Toronto / \texttt{us},
user \texttt{admin}, Wi-Fi, and SSH were injected into the boot partition so the
node came up configured headless. A genuine WPA2 validation failure
(``Password is too short — min 8 characters'') was captured and corrected.

\evfig{fig-0.6-custom-wifi-validation.png}{0.9}{WPA2-PSK minimum-length validation enforced before write.}

\subsection{First contact}
No monitor or keyboard was ever attached. The node registered via mDNS and was
reached at \texttt{dc-node-01.local}.

\evfig{fig-0.11-putty-login-admin.png}{0.62}{First SSH login as \texttt{admin} over mDNS.}

% ============================================================= PART A
\parttitle{Part A}{The Setup Manual — Execution Record}

\section{Baseline capture}\label{sec:baseline}
Before any change, the starting state was recorded — the ``before'' half of every
later claim.

\evfig{fig-01-hostnamectl-baseline.png}{0.72}{\texttt{hostnamectl} — Trixie, kernel 6.18.34, arm64. Machine ID \texttt{c8624351\ldots} (recorded; becomes rebuild proof).}

\subsection{Hardware and network}
\texttt{/proc/device-tree/model} reads \texttt{Raspberry Pi 5 Model B Rev 1.1}
directly from the device tree. \texttt{ip a} shows \texttt{eth0} as
\texttt{NO-CARRIER} and \texttt{wlan0} on a \emph{dynamic} DHCP lease
(\texttt{valid\_lft 172252sec}) — the single most important baseline finding,
and the problem §\ref{sec:static} exists to remove.

\evfig{fig-03-ip-a-baseline.png}{0.82}{\texttt{ip a} baseline: \texttt{eth0} down, \texttt{wlan0} on a dynamic lease.}

\subsection{Storage}\label{sec:storage}
\begin{lstlisting}[style=console]
Filesystem      Size  Used Avail Use% Mounted on
/dev/mmcblk0p2   58G  6.8G   49G  13% /
/dev/mmcblk0p1  505M   87M  418M  18% /boot/firmware
\end{lstlisting}
\noindent\fail\ against the manual's SSD requirement: root is on
\texttt{/dev/mmcblk0p2} — \textbf{microSD}. The ``USB Mass Storage Device'' was a
card reader holding the SD. The build continued on microSD deliberately; the
endurance limitation is carried into the risk register (§\ref{sec:security}).

\evfig{fig-04-df-h-mmcblk.png}{0.82}{\texttt{df -h} — root on \texttt{mmcblk0p2}; the manual's SSD verification fails here.}

\section{System update and firmware}
\texttt{apt full-upgrade} pulled 60 packages on a weeks-old image — the practical
argument for automatic updates (§\ref{sec:unattended}). After reboot,
\texttt{rpi-eeprom-update} reported the bootloader \textbf{up to date}
(\texttt{CURRENT = LATEST}) — the only check that means anything, since EEPROM
updates apply at boot.

\evfig{fig-18-eeprom-verified.png}{0.72}{Firmware verified up to date after reboot.}

\section{Base configuration (\texttt{raspi-config})}
The tool's title bar independently confirms \texttt{Raspberry Pi 5 Model B Rev
1.1, 8GB}. The hostname field arrived \emph{pre-populated} with \texttt{dc-node-01}
— the payoff of first-boot customisation: configuration collapsed to verification.

\evfig{fig-08-raspi-config-hostname.png}{0.62}{Hostname pre-populated; step becomes a verification, not a change.}

\section{Static addressing — failure, lockout, and rebuild}\label{sec:static}
This is the longest section because it is the only one where the build broke —
four faults in sequence, each a distinct lesson, and a Black Start recovery.

\subsection{Discover before you modify}
\texttt{nmcli connection show} revealed the active profile as
\texttt{netplan-wlan0-Calypso}, \emph{not} the manual's \texttt{"Wired connection
1"}. The \texttt{netplan-} prefix signals a renderer above NetworkManager — a
change made only with \texttt{nmcli} can be silently reverted.

\evfig{fig-10-nmcli-discovery.png}{0.9}{Real profile name discovered; Netplan-over-NetworkManager layering exposed.}

\subsection{The lockout}\label{sec:lockout}
The manual's \emph{example} values (\texttt{192.168.1.50/24}) were applied to a
node on \texttt{10.0.0.0/24}. \texttt{nmcli} accepted the syntactically valid,
semantically fatal address and returned success. The node vanished.

\evfig{fig-13-putty-connection-abort.png}{0.5}{The result: PuTTY ``Software caused connection abort.'' A lockout via the \emph{network} layer, not SSH.}

\noindent This is the report's central technical result: \textbf{a command that
exits zero is not evidence.} \texttt{nmcli} validated the address's syntax, not
its reachability. Diagnosis from the workstation showed a flood of
\texttt{169.254.x.x} (APIPA) — no DHCP was answering on that segment, separating
``the node is misconfigured'' from ``this segment has no DHCP.''

\subsection{Black Start recovery}\label{sec:rebuild}
With \texttt{eth0} unconnected and no monitor attached, no in-place repair existed.
Recovery was a full re-image using the identical Part 0 procedure — a partial
Mode-4 (Black Start) event. The proof it was a genuine rebuild:

\evfig{fig-15-hostnamectl-rebuilt.png}{0.72}{Post-rebuild \texttt{hostnamectl}: Machine ID changed \texttt{c8624351\ldots}$\rightarrow$\texttt{9e98c91d\ldots} — unforgeable evidence of a fresh install. Everything else reproduced identically.}

\subsection{The static address, applied correctly}
Post-rebuild the profile was plain \texttt{Calypso} (the \texttt{netplan-} prefix
gone — risk T5 resolved by the rebuild). Applied with this network's real values,
\texttt{ip a} confirmed the result:

\begin{center}
\begin{tabular}{@{}lll@{}}
\toprule
 & \textbf{Baseline (Fig.\ 3)} & \textbf{After}\\
\midrule
Address  & \texttt{10.0.0.249/24} & \texttt{10.0.0.2/24}\\
Flags    & \texttt{dynamic}       & \texttt{noprefixroute} (no \texttt{dynamic})\\
Lifetime & \texttt{172252sec}     & \texttt{forever}\\
\bottomrule
\end{tabular}
\end{center}
\noindent \texttt{valid\_lft forever} with no \texttt{dynamic} flag is the
definitive signature of a manual address. \pass

\subsection{The node outran its static address}\label{sec:relocate}
Relocation moved the node across three networks in three sessions — Wi-Fi
\texttt{10.0.0.2}, wired \texttt{192.168.1.2} (DHCP), and an iOS hotspot
\texttt{172.20.10.4/28}. \textbf{A host-side static address is bound to one
subnet}; the durable control for ``always reachable'' is a router-side DHCP
reservation keyed to the MAC. Recorded as a design change for the permanent
installation (and a P01 carry-forward).

\section{Users, keys, and daemon hardening}
\subsection{Account model and the permission lesson}
\texttt{groups} confirmed \texttt{sudo}, so no new admin account was needed. The
sequence \texttt{passwd} (denied) $\rightarrow$ \texttt{sudo passwd} (succeeded)
cleanly demonstrated the \texttt{/etc/shadow} permission model — an unprivileged
user may change only their own password. A deliberate \texttt{testuser1} was
created to satisfy the account-lifecycle requirement.

\subsection{Key generation and distribution}
An Ed25519 keypair was generated on the workstation and the public half pushed
with \texttt{scp} — the correct direction (a private key never leaves its origin).
The first-connection host-key fingerprint prompt is a security control, not a
formality.

\evfig{fig-30-key-login-no-password.png}{0.72}{Key-based login with no password prompt — pubkey auth working end to end, verified in a second session while the first stayed open.}

\subsection{Validate-before-replace}
Editing \texttt{sshd\_config} through a validated temporary copy caught a real
syntax error at line 117; the abort left the live file untouched. On a headless
node an \texttt{sshd} that will not start is an unrecoverable lockout — this habit
is what stands between a typo and a trip to the hardware.

\section{Host firewall — \texttt{ufw}}
Default-deny inbound, allow outbound, SSH explicitly permitted — and crucially
\texttt{ufw allow OpenSSH} was issued \emph{before} \texttt{ufw enable}, so the
management path was never severed. Both IPv4 and IPv6 rules were created, which
matters on a node holding global IPv6 addresses.

\evfig{fig-33-ufw-status-verbose.png}{0.82}{\texttt{ufw status verbose} — default-deny inbound, OpenSSH allowed on v4 and v6.}

\section{Brute-force protection — \texttt{fail2ban}}
The \texttt{sshd} jail is active with the \textbf{journal} backend
(\texttt{\_SYSTEMD\_UNIT=ssh.service}) — the modern source; a file-based filter on
a journal-only system would watch a file that never fills.
\texttt{jail.conf $\rightarrow$ jail.local} is the upgrade-safe customisation pattern.

\evfig{fig-35-fail2ban-jail-sshd.png}{0.82}{\texttt{fail2ban-client status sshd} — active jail, journal backend, zero bans (clean baseline).}

\section{Automatic security updates}\label{sec:unattended}
\texttt{unattended-upgrades} installed and \textbf{enabled at boot} (systemd
symlink into \texttt{multi-user.target.wants}). This answers the 60-pending-package
observation: a node left alone drifts out of patch currency within days.

\section{Time synchronisation}
\texttt{timedatectl} showed \texttt{synchronized: yes}, \texttt{NTP active}, and the
correct \texttt{-0400 EDT} offset with \texttt{RTC in local TZ: no} (hardware clock
in UTC — the correct posture). Correct time is load-bearing for \texttt{fail2ban}
correlation, TLS validity, and any future TOTP.

\evfig{fig-38-timedatectl-synchronized.png}{0.82}{\texttt{timedatectl} — clock synchronised, NTP active.}

\section{Core tooling and portfolio workflow}
Tooling installed (\texttt{git} 2.47.3, \texttt{vim}, \texttt{tmux}, \ldots), with
\texttt{dnsutils}$\rightarrow$\texttt{bind9-dnsutils} shown as a virtual-package
resolution rather than an error. GitHub SSH auth was proven — the pre-registration
\texttt{Permission denied (publickey)} is expected, and the success greets the
account by name.

\evfig{fig-41-github-auth-success.png}{0.72}{\texttt{Hi RedjiJB!} — key bound to the correct GitHub identity.}

\section{Container substrate — Docker}
The install script's echoed actions show a GPG key placed in
\texttt{/etc/apt/keyrings} and a \texttt{signed-by} repository line, so packages
are verified through APT's normal trust path. \texttt{hello-world} pulled the
\texttt{arm64v8} variant — multi-arch resolution working, vindicating the 64-bit
OS choice.

\evfig{fig-43-docker-hello-world.png}{0.72}{Docker daemon functional; \texttt{arm64v8} image resolution.}

% ============================================================= PART B
\parttitle{Part B}{Meta Report}

\section{Objective and scope}
Provision a Raspberry Pi 5 from bare hardware into a hardened, reproducible
server node — the foundation for the D-Central and CST portfolios and the P01
Black Start bootstrap. In scope: image prep, baseline, patching, addressing, key
hardening, firewall, intrusion protection, auto-patching, time sync, tooling,
portfolio workflow, container substrate. Out of scope: service deployments
(later projects).

\section{Competencies demonstrated}
Linux administration (accounts, the \texttt{/etc/shadow} model, systemd, package
management); network configuration and diagnosis (lease semantics, NetworkManager/
Netplan layering, APIPA identification, mDNS, NAT64 observation); host hardening
(Ed25519 keys, host-key verification, dual-stack default-deny, journal-backed
\texttt{fail2ban}, validate-before-replace); and service management with an
honest, evidence-backed record of a self-inflicted outage and its recovery.

\section{Analysis and discussion}
\textbf{First-boot customisation is the highest-leverage step} — it turned
configuration into verification and made the rebuild cheap. \textbf{The central
finding} is that \texttt{nmcli} succeeded and the node disappeared: exit status
measured the parser, not the outcome, which is the strongest possible case for
the \pass-discipline used throughout. \textbf{Reproducibility was proven by
recovery, not assertion} — the Machine ID change documents a rebuild whose cost
was one imaging pass. \textbf{Static addressing solved the wrong half of the
problem} for a portable node; the durable control is a router-side reservation.

\section{Challenges and troubleshooting}
\subsection{Bookworm $\rightarrow$ Trixie}\label{sec:trixie}
The manual targets Debian 12; the recommended image is Debian 13, with Bookworm
demoted to \emph{Legacy}. Trixie was accepted deliberately. The cost: the manual
can no longer be trusted verbatim, and §\ref{sec:static} is where that bit. A
build manual pinned to a distribution release has a shelf life.

\subsection{The include-ordering trap (D3)}\label{sec:d3fix}
The hardening drop-in existed, validated, and loaded — yet
\texttt{PasswordAuthentication} stayed \texttt{yes}. \texttt{sshd\_config}
resolves the \emph{first} value obtained, and cloud-init's
\texttt{50-cloud-init.conf} sorted ahead of the \texttt{99-} file. Renaming to
\texttt{00-} won the ordering. \textbf{Configuration that is present is not
configuration that is effective.}

\evfig{fig-59-negative-test-permission-denied.png}{0.82}{The negative test — \texttt{Permission denied (publickey)} — proves password auth is \emph{refused}, the evidence the whole hardening objective rests on.}

\subsection{One stalled activation, three symptoms (D1 + D8)}
A single stalled DHCP activation produced a dead resolver, an unsynchronised clock,
and dropping SSH sessions — none of which names the network on its own. Moving the
node to a static Wi-Fi profile fixed all three at once. \texttt{NTP service:
active} described a daemon synchronising with nothing — the same present-vs-effective
lesson as D3, in a second subsystem.

\section{Security and risk considerations}\label{sec:security}
Controls in place and evidenced: non-root \texttt{admin} with sudo; Ed25519 key
auth with host-key verification; \textbf{key-only SSH proven by negative test};
dual-stack default-deny firewall; active \texttt{fail2ban} jail; auto-patching
enabled at boot; NTP synchronised; APT GPG trust for Docker; per-machine keypairs.
Residual risks: physical access; microSD endurance; addressing instability on a
portable node; \texttt{docker}-group root-equivalence (an accepted, documented
trade-off); mDNS advertisement on the local segment.

\section{Conclusion and lessons learned}
A hardened, documented, container-ready node exists as the portfolio and D-Central
foundation, with all nine defects closed and reboot persistence verified. Ordered
by what they cost to learn:
\begin{enumerate}[leftmargin=1.4em]
\item A command that exits zero is not evidence — verify the outcome, not the return code.
\item Configuration that is \emph{present} is not configuration that is \emph{effective}.
\item Never copy example values; every address is a placeholder until checked against \texttt{ip a}.
\item Enumerate every path that can sever your own management channel.
\item Reproducibility is proven by recovery, not by assertion.
\item Discover before you modify — the connection name changed three times in one build.
\item Watch for configuration layered above the tool you are using.
\item Validate before replacing.
\item Static addressing solves the wrong half of the problem for a portable node.
\item A build manual pinned to a distribution release expires.
\end{enumerate}

% ============================================================= §Defects
\section{Defect Register}
\begin{longtable}{@{}c p{6.6cm} c@{}}
\toprule
\textbf{ID} & \textbf{Defect} & \textbf{Status}\\
\midrule\endhead
D1 & DNS secondary typo \texttt{8.4.4.4} $\rightarrow$ \texttt{8.8.4.4} & \closed\\
D2 & Pre-existing \texttt{redji} account with sudo & \closed\ (absent on current install)\\
D3 & Password/root login enabled — include-ordering & \closed\ (negative test)\\
D4 & Docker required \texttt{sudo} & \closed\\
D5 & \texttt{\~{}/.ssh} permissions unverified & \closed\\
D6 & \texttt{testuser1} held sudo & \closed\ (revoked)\\
D7 & Reboot survival unverified & \closed\ (6 controls persisted)\\
D8 & Clock unsynchronised after reboot & \closed\ (root cause: D1 chain)\\
D9 & Portfolio push & \closed\ (published, sanitised)\\
\bottomrule
\end{longtable}

% ============================================================= §Rubric
\section{Rubric Self-Assessment}
\begin{longtable}{@{}p{5.6cm} c p{4.4cm}@{}}
\toprule
\textbf{Dimension} & \textbf{/4} & \textbf{Note}\\
\midrule\endhead
Objective clarity          & 4 & Scope and exclusions stated\\
Reproducibility            & 4 & Proven by rebuild under adversarial conditions\\
Evidence quality           & 4 & 62 figures; failed verifications retained\\
Technical correctness      & 3 & Deviations analysed, not smoothed\\
Analysis depth             & 4 & Mechanism explained, not just outcome\\
Professional communication & 4 & Defect register, honest failure reporting\\
Competency mapping         & 4 & Four courses, figure-level backing\\
Security consideration     & 4 & Audited, failed, fixed, proven by negative test\\
\midrule
\textbf{Total} & \textbf{32/32} & Above threshold ($\geq$28)\\
\bottomrule
\end{longtable}

% ============================================================= Appendix
\parttitle{Appendix C}{Evidence Index and File Mapping}

Rename screenshots to these names in \texttt{evidence/} so figure references
resolve and the \verb|\evfig| macro finds each file.

\begin{longtable}{@{}c p{9.6cm}@{}}
\toprule
\textbf{Fig} & \textbf{Filename}\\
\midrule\endhead
0.1--0.12 & \texttt{fig-0.1-imager-device-pi5.png} \ldots\ \texttt{fig-0.12-imager-write-complete.png}\\
1  & \texttt{fig-01-hostnamectl-baseline.png}\\
2  & \texttt{fig-02-device-tree-model.png}\\
3  & \texttt{fig-03-ip-a-baseline.png}\\
4  & \texttt{fig-04-df-h-mmcblk.png}\\
5  & \texttt{fig-05-apt-full-upgrade.png}\\
6--9 & \texttt{fig-06-raspi-config-main.png} \ldots\ \texttt{fig-09-raspi-config-localisation.png}\\
10 & \texttt{fig-10-nmcli-discovery.png}\\
11 & \texttt{fig-11-nmcli-quote-error.png}\\
12 & \texttt{fig-12-nmcli-wrong-subnet.png}\\
13 & \texttt{fig-13-putty-connection-abort.png}\\
14 & \texttt{fig-14-powershell-apipa.png}\\
15 & \texttt{fig-15-hostnamectl-rebuilt.png}\\
16--18 & \texttt{fig-16-supd-typo-52-packages.png} \ldots\ \texttt{fig-18-eeprom-verified.png}\\
19--24 & \texttt{fig-19-nmcli-calypso-renamed.png} \ldots\ \texttt{fig-24-df-h-reverify.png}\\
25--31 & \texttt{fig-25-groups-adduser-passwd.png} \ldots\ \texttt{fig-31-sshd-config-syntax-error.png}\\
32--39 & \texttt{fig-32-ufw-install.png} \ldots\ \texttt{fig-39-core-tooling.png}\\
40--44 & \texttt{fig-40-ssh-keygen-github-denied.png} \ldots\ \texttt{fig-44-docker-compose-version.png}\\
45--52 & \texttt{fig-45-sshd-T-audit-failed.png} \ldots\ \texttt{fig-52-testuser1-has-sudo.png}\\
53--60 & \texttt{fig-53-post-reboot-part1.png} \ldots\ \texttt{fig-60-dns-activation-fail-testuser1-ntp.png}\\
61--62 & \texttt{fig-61-wired-stuck-activating-linklocal-dns.png}, \texttt{fig-62-wifi-static-dns-ntp-synchronized.png}\\
\bottomrule
\end{longtable}

\vspace{1em}
\noindent\textbf{Excluded from evidence (never publish):} online-banking capture
(account balances/transactions) and two unrelated CCNA/ROAS topology images.
Absence in the published tree was verified post-push with
\texttt{git ls-files | grep -E "071704|184342|183550"} returning empty.

\end{document}
